\scanpage{86}

\margin{\inserted{Bien entendu, en remplaçant au besoin $\varphi$ par \deleted{\unsure{$f$}} son symétrique,}}\unsure{En th. 1, 1°, hypothèse b)}

\deleted{Bien entendu, \unsure{en vertu de} l'inégalité \unsure{indiquée|analogue}, $(4)$ $\sum\varphi(-\alpha_1,\ldots,-\alpha_n)\leq\sum\varphi(-\beta_1,\ldots,-\beta_n)$}

\uline{Remarque 1} Notons \deleted{\unsure{d'ailleurs|toutefois}} que si deux suites $(\alpha_i)_{i\leq n},(\beta_i)_{i\leq n}$ satisfont aux conditions générales du \deleted{lemme}\inserted{th. 1, 1°} \inserted{$(\alpha_1\geq\ldots\geq\alpha_n,\ \alpha_1+\cdots+\alpha_{\unsure{M|N|m}}\leq\beta_1+\cdots+\beta_M\ \deleted{\unsure{et}}\ 1\leq M\leq n$} alors \deleted{il en est de les \unsure{mêmes suites} $(\alpha_i)_{i\leq p},(\beta_i)_{i\leq p}$ pour tout $1\leq p\leq n$, \unsure{(ceci pour $p=n$ bien que} $\sum_1^p\alpha_i=\sum_1^p\beta_i$), donc on peut leur appliquer la \unsure{formule} du lemme, et en particulier par exemple $\varphi(\alpha_1,\ldots,\alpha_p)$ qu'on obtient en remplaçant $\alpha_{p+1},\ldots,\alpha_n$ et $\beta_{p+1},\ldots,\beta_n$ par}

\inserted{les suites à $k$ éléments} $(\alpha_i)_{i\leq k}$ et $(\beta_i)_{i\leq k}$ \inserted{satisfont aux conditions analogues} \deleted{quelque que soit}\inserted{quel que soit} $1\leq k\leq n$. On peut donc appliquer à ces suites le même lemme, et on trouve \deleted{\unsure{donc}}\inserted{p. ex.} que \deleted{\unsure{(quitte à} \illegible{2} \unsure{le lemme}} \inserted{si $\varphi$ est une fonction convexe} croissante de \uline{$k$} arguments, on a aussi $\varphi(\alpha_1,\ldots\alpha_k)\leq\varphi(\beta_1,\ldots,\beta_k)$. La remarque analogue vaut pour les \deleted{corollaire 1 du lemme}\inserted{conditions du th. 1, 2°}.

% slanted
\margin{\uline{Remarque 2} Si on suppose $(\alpha_i)$ $(\beta_i)$ \unsure{comme} dans l'\unsure{exposé a|espace}, et que $f(\alpha_1\ldots\alpha_n)\leq f(\beta_1\ldots\beta_n)$ pour \deleted{\unsure{toute}} $f$ telle \unsure{que} \ldots{} alors on a $\sum_1^m\alpha_i\leq\sum_1^m\beta_i$ (évidemment, cela revenant à $\sup_H\sum_{i\in H}\alpha_i$, \deleted{\unsure{où}} $H$ à \deleted{\unsure{au plus}} $m$ éléments)}

\deleted{\uline{Réciproque et exemples de convexité du corollaire 1.}}

Donnons maintenant des cas intéressants de fonctions $f$ satisfaisant aux conditions du \deleted{corollaire 1}\inserted{th. 1, 2°}. \deleted{\unsure{qui est}}

Si \uline{\deleted{\unsure{$f(\rho_1$}} $f(\rho_1,\ldots\rho_n)$ est une fonction convexe} croissante des \inserted{l'ens. des} arguments $\rho_i\geq0$, alors $f(e^{t_1},\ldots,e^{t_n})$ est \deleted{\unsure{convexe}}\inserted{une} fonction convexe de \deleted{\unsure{chacun des arguments}} $(t_i)\in R^{n}$. Comme \deleted{la fonction convexe $f$ \unsure{sera celle qui est}} \inserted{$f$ sera} une enveloppe supérieure de fonctions \deleted{\unsure{linéaires}}\inserted{linéaires} affines $\inserted{\sum a_i x_i+b}$ \inserted{avec la propriété $a_i\geq0$ (!)}, $f(e^{t_1},\ldots,e^{t_n})$ sera enveloppe supérieure \deleted{et que toute fonction convexe} de fonctions $\sum a_i e^{t_i}+b$, qui sont convexes (car \deleted{\unsure{si}} les fonctions $e^{t_i}$ sont convexes, et les $a_i$ positifs), donc elle sera elle-même convexe.

En particulier, \deleted{si $f$ par une norme \deleted{fonction} de la. Un} \unsure{appelons|appl.ons}, \uline{normes de Schatten} sur $R^n$ toute \uline{norme} $N(x)=N(x_1,\ldots x_n)$ qui est symétrique, telle que $N(x_1,\ldots x_n)=N(|x_1|,\ldots|x_n|)$, \deleted{et qui est fonction croissante de chacun des arguments pour} \inserted{$x_1\geq0,\ldots x_n\geq0$}, \deleted{\unsure{et par le fait est donc} \illegible{2}} \inserted{normalisée par $N(e_1)=1$.}

% margin note written sideways
\margin{Alors on prouve que $N$ est fonction croissante \deleted{\unsure{des arguments}} \unsure{(ens)} des \unsure{longueurs} des éléments positifs ((car la boule unité, étant stable par multiplication par une \unsure{homothétie|fonction} $(\lambda_i)$ avec $|\lambda_i|=1$ pour \deleted{tout} $i$, le sera pour \deleted{\unsure{toute}} fonction majorée par $1$ en module).}

% archive sheet number: 87
